\lecture{5}{8. September 2025}{Rigid Body Equilibrium, Pt. 1}

\section{Equilibrium of a Rigid Body}
\subsection{Conditions for Rigid-Body Equilibrium}
As mentioned previously the force and couple moments acting on a body can be reduced to an equivalent resultant force and resultant couple moment at any arbitrary point $O$ on or off the body. If these resultants are both equal to zero the body is said to be in equilibrium. Mathematically this is:
\begin{align*}
  \textbf{F}_R &= \sum \textbf{F} = \textbf{0} \\
  \left( \textbf{M}_R \right)_= &= \sum \textbf{M}_O = \textbf{0}
.\end{align*}
These two equations are both necessary and sufficient for equilibrium. This is true in both 2 and 3 dimensions.

\subsection{Characteristics of Dry Friction}
Friction is the name of the force that resists the movement between two contacting surfaces that slide relative to each other. Here only dry friction, also called \textit{Coulomb friction}, is described. This occurs when there is no lubricating liquid between the two contacting surfaces.

This force will resist motion caused by an applied force $\textbf{P}$ when the size of the applied force $\textbf{P}$ increases above the \textit{limiting static frictional force}, $F_s$, motion will occur. $F_s$ is given by:
\[ 
F_s = \mu_s N
\]
where $\mu_s$ is the coefficient of static friction and $N$ is the normal force. When the applied force $\textbf{P}$ increases above $F_s$ and motion occurs the friction force will drop to a smaller value called the \textit{kinetic frictional force} given by:
\[ 
F_k = \mu_k N
\]
where $\mu_k$ is the coefficient of kinetic friction.
